% Generated draft: requires mathematical and source review.
\documentclass[11pt]{article}
\usepackage[margin=25mm]{geometry}
\usepackage{fontspec}
\setmainfont{Latin Modern Roman}
\usepackage{amsmath,amssymb,mathtools,unicode-math}
\setmathfont{Latin Modern Math}
\usepackage{enumitem,needspace,xurl,hyperref}
\setlength{\parindent}{0pt}
\setlength{\parskip}{5pt}
\setlength{\emergencystretch}{3em}
\newcommand{\N}{\mathbb{N}}
\newcommand{\Z}{\mathbb{Z}}
\newcommand{\Q}{\mathbb{Q}}
\newcommand{\R}{\mathbb{R}}
\newcommand{\C}{\mathbb{C}}
\newcommand{\F}{\mathbb{F}}
\newcommand{\A}{\mathbb{A}}
\newcommand{\PP}{\mathbb{P}}
\newcommand{\E}{\mathbb{E}}
\newcommand{\eps}{\varepsilon}
\newcommand{\dd}{\,\mathrm{d}}
\begin{document}
\section*{UnsolvedMath ID: 30000556}
\textbf{Source problem code:} catalogue entry
\subsection{Definitions and mathematical statement}
For the standard closed-manifold setting, let $(M,\omega)$ be a closed symplectic manifold and let $\operatorname{Ham}(M,\omega)$ be its group of Hamiltonian diffeomorphisms. A Hamiltonian path $\phi_H^t$, generated by $H_t:M\to\R$, has Hofer length
\[
\ell_H(\phi_H)=\int_0^1\left(\max_M H_t-\min_M H_t\right)\,\dd t.
\]
The Hofer distance is the infimum of this length over Hamiltonian paths with prescribed endpoints. Addition of a time-dependent constant to $H$ does not change the length.

The question asks whether every
\[
a\in\pi_1(\operatorname{Ham}(M,\omega),\operatorname{id})
\]
has a representative $\gamma$ that is a closed geodesic for the Hofer metric. Thus $\gamma$ is a loop based at the identity with $[\gamma]=a$, satisfying the intended local geodesic condition, including across its closing point. A metric-local interpretation requires sufficiently short subarcs of $\gamma$ to realize the Hofer distance between their endpoints. A closed geodesic is not required to minimize length globally among all loops or all paths with its coinciding endpoints. For the identity fundamental-group element, a constant loop is allowed unless the source requires nonconstant geodesics.

The extract does not identify $(M,\omega)$, specify whether all symplectic manifolds or a particular class is intended, or define the group topology and geodesic regularity/local-minimality convention. The usual $C^\infty$ topology on the Hamiltonian group and the closed-manifold setting above are explicit contextual choices. Open manifolds require support conventions, and variational versus metric meanings of Hofer geodesic require distinction; neither is silently imposed as a verified source hypothesis.
\subsection{Short English statement}
Does every fundamental-group class of a Hamiltonian diffeomorphism group contain a closed Hofer geodesic? The manifold class, topology and precise geodesic convention remain unspecified; the target concerns local geodesic representatives rather than globally shortest loops.
\subsection{Sources}
\begin{itemize}
\item Supplied report: \emph{Geometric Group Theory, Hyperbolic Dynamics and Symplectic Geometry} (2007), Oberwolfach Reports, DOI 10.4171/owr/2006/33, pp. 1991--2058; proposed year 2006. \url{https://doi.org/10.4171/owr/2006/33}.
\item Local UnsolvedMath registry, permanent ID 30000556, catalogue entry; created 25 August 2026, updated 16 September 2026. The statement is labelled exact and confirmed without textual change. The supplied 21 August 2026 literature triage labels it open and reports no authoritative resolution of the extracted claim. No further external reference is supplied, and no independent status check is asserted here.
\end{itemize}
\end{document}
